% =====================================================================
% Stage27 Manuscript Integration
% Generated only from frozen Stage27-4A artifacts.
% =====================================================================

\subsection{Unseen Attack-Family Generalization}

\subsubsection{Chronology-first audit design}

Stage27 evaluated zero-training-exposure attack-family generalization under a strict \texttt{TRAIN < VALIDATION < TARGET} protocol. The held-out family was absent from both training and validation, thresholds were selected only on known-family validation data, and the target was not used for model selection, threshold tuning, or adaptation. The experiment is therefore described as an unseen attack-family generalization audit rather than as formal proof of zero-day detection.

Seven primary families were preregistered. Five were structurally executable. DOS was ineligible because the available earlier day-atomic training period contained no known-family attack positives, whereas AUTH\_BRUTE\_FORCE was ineligible because insufficient earlier weekday depth existed for separate training and validation periods. INFILTRATION was executable but is reported descriptively only because the held-out support was 36.

\subsubsection{Primary unseen-family ranking}

Unseen-family ranking was strongly family dependent. DDoS retained near-perfect discrimination: XGBoost reached ROC-AUC 0.9982 and PR-AUC 0.9925, while LightGBM reached ROC-AUC 0.9986 and PR-AUC 0.9940. Web Attack also transferred strongly, with XGBoost ROC-AUC 0.9693 and PR-AUC 0.7206 and LightGBM ROC-AUC 0.9901 and PR-AUC 0.7605.

Bot traffic showed substantial collapse. XGBoost produced ROC-AUC 0.3224 and PR-AUC 0.003256, while LightGBM produced ROC-AUC 0.5591 and PR-AUC 0.004879. The XGBoost PR-AUC was below the target prevalence anchor, giving negative PR-excess -0.001467. Port Scan was materially learner-dependent: XGBoost reached ROC-AUC 0.5506 compared with 0.7559 for LightGBM.

The overall Stage27 outcome is therefore selective family transfer rather than universal unseen-family generalization.

\subsubsection{Frozen operating-point transfer}

Threshold-independent ranking did not guarantee useful frozen-threshold detection. At the BALANCED operating point, DDoS recall was 66.20\% for XGBoost and 26.25\% for LightGBM despite ROC-AUC above 0.998 for both learners. Similarly, LightGBM retained Port Scan ROC-AUC 0.7559 while BALANCED recall was only 1.17\%.

These results distinguish ranking generalization from operating-point transfer and support the frozen Stage27 outcome of ranking--threshold divergence.

\subsubsection{Behavioral similarity}

The secondary behavioral-similarity audit did not show a monotonic relationship with transfer. BOT had the highest observed similarity to a seen family (0.6994) but weak generalization, whereas DDoS had lower similarity (0.3834) and near-perfect ranking. Behavioral proximity under the frozen 11-descriptor centroid definition therefore does not appear sufficient by itself to explain the transfer pattern.

\subsection{Discussion of Attack-Family Novelty}

Stage27 shows that strong performance on known attack families cannot be interpreted as evidence of uniform robustness to attack-family novelty. DDoS and Web Attack retained strong ranking for both learners, Bot collapsed, and Port Scan exhibited substantial learner dependence.

The experiment also reinforces a broader finding across the study: threshold-independent discrimination and fixed operating-point behavior are distinct properties. Representation-specific chronological evaluation, the temporal-validation stress test, cross-dataset transfer, and now unseen-family transfer each expose cases in which ranking and frozen-threshold behavior diverge. Reporting only ROC-AUC or PR-AUC would therefore provide an incomplete description of deployment robustness.

The evidence does not establish a universal learner winner. XGBoost and LightGBM agree closely on DDoS and Web Attack ranking but differ substantially for Port Scan and Bot. Learner dependence is therefore itself family-dependent.

\subsection{Stage27 Limitations}

\begin{itemize}
\item Only five of seven preregistered families were structurally executable under strict chronology.
\item INFILTRATION is descriptive only because the held-out support was 36.
\item The design is chronology-first zero-training-exposure evaluation rather than textbook LOAO in which every other family is necessarily represented in training.
\item The 95\% bootstrap intervals quantify target-sampling uncertainty conditional on the fitted model and do not include retraining, seed, model-selection, or broader population uncertainty.
\item Behavioral similarity uses only 11 preregistered aggregate descriptors and is descriptive only.
\item The benchmark-specific results do not establish universal real-world zero-day detection.
\end{itemize}

% Main Stage27 figures:
% stage27_primary_roc_auc_ci.png
% stage27_balanced_recall_ci.png
%
% Supplementary:
% stage27_primary_pr_auc_ci.png
